% Adapted from topics/oop/easy/adt.tex for dataclass-based compound data.
\begin{blocksection}
\question \begin{parts}
\part What is an \define{Abstract Data Type}?
\begin{solution}[0.65in]
An abstract data type describes the information a value represents and the
operations we can use with it, without requiring users to know how those
operations are implemented. For example, a session record can provide a topic,
room, and duration through named attributes.
\end{solution}
\part What is the relationship between a class and an ADT?
\begin{solution}[0.65in]
A class is one way to implement an ADT. A dataclass can bundle a session's
information into fields such as \lstinline{topic}, \lstinline{room}, and
\lstinline{minutes}. Calling \lstinline{Session("Recursion", 310)} constructs
an instance, and reading \lstinline{s.room} accesses its room. Code using this
interface need not know how the fields are stored. The class's definition
implements the representation; the meaning and supported operations form the
abstraction.
\end{solution}
\end{parts}
\end{blocksection}
\begin{questionmeta}
\textbf{Teaching Tips}\par\nopagebreak[4]
Suggested Time: 3--4 min; Difficulty: Conceptual
\nopagebreak[4]
\begin{itemize}
    \item \textbf{Connect to functions:} Using a function without inspecting its body is analogous to using a compound value through its interface.
    \item \textbf{Make the distinction concrete:} A dataclass provides a convenient implementation; the abstraction is the meaning of the data and what its interface lets us do. Use the session-record example in the solution.
\end{itemize}
\end{questionmeta}
